\documentclass[10pt,a4paper]{amsart}
\usepackage[margin=19mm]{geometry}
\usepackage{amsmath,amssymb,mathtools}
\usepackage[scr=rsfs,cal=euler]{mathalfa}
\usepackage{xfrac}
\usepackage[unicode,hidelinks,bookmarks=false]{hyperref}
\newcommand{\ie}{i.e.\ }
\newcommand{\cf}{cf.\ }
\newcommand{\resp}{respectively\ }
\newcommand{\Subsection}{\subsection}
\providecommand{\nobreakdash}{\nobreak\hbox{-}}
\setlength{\emergencystretch}{2em}
\multlinegap0pt
\usepackage{bm}
\usepackage{enumerate}
\usepackage{amssymb}
\usepackage{booktabs}
\usepackage{float}
\usepackage{multirow}
\usepackage{algorithm}\usepackage{algpseudocode}
\usepackage{dsfont}
\usepackage{caption}
\usepackage{tikz}
\newtheorem{thm}{Theorem}
\newtheorem{lem}{Lemma}
\title{{ Error Analysis of Parameter Prediction via Gaussian Process Regression and Its Application to Weighted Jacobi Iteration}}
\author{Tiantian Sun, Juan Zhang}
\date{Source: arXiv:2602.22679v1 (CC BY 4.0)}
\begin{document}
\maketitle

\noindent\emph{Source and licence.} { Error Analysis of Parameter Prediction via Gaussian Process Regression and Its Application to Weighted Jacobi Iteration}. Version arXiv:2602.22679v1 (CC BY 4.0), distributed by arXiv under the Creative Commons Attribution 4.0 International licence, http://creativecommons.org/licenses/by/4.0/, as recorded in the arXiv licence metadata for this exact version and shown on the abstract page https://arxiv.org/abs/2602.22679v1. Full licence notice: \texttt{licenses/arxiv-2602.22679-CC-BY-4.0.txt}.\par\smallskip

\noindent\textit{Editorial scope.} Counted unit: the article's thm thm\_GPR\_nf (source0.tex lines 317--324). The article's complete original proof of it is delivered with it (source0.tex lines 325--374, one \texttt{proof} environment of 50 lines). No mathematics is written, repaired or re-derived: the statement and the proof are the article's own lines, in the article's own notation, and no retained line is substituted or re-flowed.  Retained uncounted as the source-local context the proof needs: lines 256--266; lines 285--294.  Nothing was omitted from any retained span, so the package carries no omission record.  No retained line contains a citation, so the wrapper carries no bibliography and no bibliography entry.  No retained line contains a figure environment, an image-inclusion command or drawing code, so no image is delivered and the package declares no asset dependency.  Editorial formatting only, in addition: an AMS \texttt{amsart} preamble that restates the article's own declarations and macros used by the retained lines, namely the environments \texttt{thm}, \texttt{lem} on separate counters, together with the standard packages those lines need.  The retained environments print in this document as Theorem 1, Lemma 1 in order of appearance, whereas the article prints the same items with the numbering of its own full document.  The complete native source of the article as distributed in the arXiv e-print for this exact version is bundled unchanged as \texttt{sources/195345/source0.tex}.
\begin{thm}
	\label{thm_Xi}
	Assume $f(\boldsymbol{z}) = \sum\limits_{j=1}^{\infty} \beta_j \kappa (z,x_j) \in \mathcal{N}_\mathcal{H}$, where $\kappa(z,x)$ is a reproducing kernel associated with a Hilbert space $\mathcal{H}$. Let  $p_1,p_2,\dots,p_q$ be a set of linearly independent functions in $C(\Omega)$. 
	%(forming a basis of a $q$-dimensional subspace ). 
	For a subset $\Xi = \{\xi_1,\xi_2,\dots,\xi_q \} \subseteq \Omega$, the function $f(z)$ can be rewritten as:
	\begin{equation}
		\label{pro_f_z}
		f(\boldsymbol{z}) = \sum\limits_{i = 1}^q f(\xi_i) p_i(\boldsymbol{z}) +\left\langle f,\kappa(\cdot,\boldsymbol{z})-\sum\limits_{i= 1}^q p_i(\boldsymbol{z})\kappa(\cdot,\xi_i) \right\rangle_\mathcal{H}.
	\end{equation}
	
\end{thm}

\begin{lem}
	\label{lem_g}
	Let $\kappa(x, x_1), \kappa(x, x_2), \dots, \kappa(x, x_N)$ be a set of linearly independent functions, and let $p_1(x), \dots, p_q(x)$ be another set of linearly independent functions. Let $A = (a_{ij})$ be an invertible matrix, and define the functions $g_i(x)$ as:
	$$
	g_i(x) = \sum_{j=1}^N \sum_{l=1}^N \kappa(x, x_j) a_{jl} p_i(x), \quad i = 1, \dots, q,
	$$
	then, $g_1(x), g_2(x), \dots, g_q(x)$ are also linearly independent functions.
	
	
\end{lem}

\begin{thm}
	\label{thm_GPR_nf}
	Let $\Omega \subseteq \mathbb{R}^d$ be open. Suppose that $\kappa$ is a positive definite kernel on $\Omega$. Then for every $\boldsymbol{z}\in \Omega$ the error between true function $f\in \mathcal{N}_\mathcal{H} (\Omega)$ and its prediction $f^*$ from GPR can be bounded by 
	\begin{equation}
		\vert f(\boldsymbol{z})-f^*(\boldsymbol{z})\vert^2 \le C\cdot Var(f^*(\boldsymbol{z})),
	\end{equation}
	where $C$ is an nonnegative constants, and $Var(f^*(\boldsymbol{z}))$ is the variance of $f^*(\boldsymbol{z})$.
\end{thm}

\begin{proof}
	Let $p_1(x), p_2(x),\dots,p_q(x)$ be a set of linearly independent functions, $\mathcal{P}$ is a space spanned by $\{p_1,p_2,\dots,p_q\}$, and choose a subset $\Xi = \{\xi_1,\xi_2,\dots,\xi_q \} \subseteq \Omega$ . Assume $\dim(\mathcal{P})=q \le N$, and let $\{p_i(\boldsymbol{z})\}_{i = 1}^q$ is a basis of $\mathcal{P}$.
	%, which can be reproduced by $K$, i.e.  $p_i(\boldsymbol{x}) = \sum\limits_{j = 1}^N \kappa(\boldsymbol{z},\boldsymbol{x}_j)p_i(\boldsymbol{x}_j)$. 
	Then, according to Theorem \ref{thm_Xi}, $f(\boldsymbol{z})$ can be rewritten as 
	\begin{equation}
		\label{f_z}
		f(\boldsymbol{z}) = \sum\limits_{i = 1}^q f(\xi_i) p_i(\boldsymbol{z}) +\left\langle f,\kappa(\cdot,\boldsymbol{z})-\sum\limits_{i= 1}^q p_i(\boldsymbol{z})\kappa(\cdot,\xi_i) \right\rangle_\mathcal{H}.
	\end{equation}
	
	
	
	
	Let $g_i(x) :=\sum\limits_{j =1}^N\sum\limits_{k=1}^N \kappa(x,x_j)a_{jk} p_i(x_k)   = \sum\limits_{j=1}^N \kappa(x,x_j) g_i(x_j)$, where $g_i(x_j) = \sum\limits_{k = 1}^N a_{jk}p_i(x_k)$. 
	Let $U^*(\boldsymbol{z}) =  K_X(\boldsymbol{z}) K_{XX}^{-1}$, then $\bar{f^*}(\boldsymbol{z})  = U^*(\boldsymbol{z})Y$. {  There is}
	
	
	
	$$
	\begin{aligned}
		f^*(\boldsymbol{z}) &= U^*(\boldsymbol{z})Y =\sum\limits_{j = 1}^N u_j^*(\boldsymbol{z})f(x_j)\\
		&=\sum\limits_{j = 1}^N u_j^*(\boldsymbol{z}) \left(\sum\limits_{i = 1}^q f(\xi_i) p_i(x_j) +\left\langle f,\kappa(\cdot,x_j)-\sum\limits_{i= 1}^q p_i(x_j)\kappa(\cdot,\xi_i) \right\rangle_\mathcal{H} \right)\\
		&=\sum\limits_{i = 1}^q \sum\limits_{j = 1}^N u_j^*(\boldsymbol{z}) f(\xi_i) p_i(x_j)+ \left\langle f, \sum\limits_{j = 1}^N u_j^*(\boldsymbol{z})\kappa(\cdot,x_j)- \sum\limits_{i= 1}^q\sum\limits_{j = 1}^N u_j^*(\boldsymbol{z})p_i(x_j)\kappa(\cdot,\xi_i) \right\rangle_\mathcal{H}.
	\end{aligned}
	$$
	
	Since $U^*(\boldsymbol{z}) =\begin{bmatrix}u_1^*(\boldsymbol{z}) & u_2^*(\boldsymbol{z}) &\cdots &u_N^*(\boldsymbol{z}) \end{bmatrix}= K_X (\boldsymbol{z}) (K_{XX})^{-1}$,  $\sum\limits_{i = 1}^q \sum\limits_{j = 1}^N u_j^*(\boldsymbol{z}) f(\xi_i) p_i(x_j) $ can be rewritten as $ K_X(\boldsymbol{z}) K_{XX}^{-1} P^T Y_{\xi}$, where $P = (p_i(x_j))\in \mathbb{R}^{q\times N}$, $Y_{\xi} =(f(\xi_1),f(\xi_2), \dots, f(\xi_q) )^T $. Let $(g_1(\boldsymbol{z}),\dots,g_q(\boldsymbol{z})) =K_X(\boldsymbol{z}) K_{XX}^{-1} P^T $, according to Lemma \ref{lem_g}, $g_1(\boldsymbol{z}), g_2(\boldsymbol{z}),\dots,g_q(\boldsymbol{z})$ are linearly independent. Obviously, we have 
	$$
	\sum\limits_{i = 1}^q \sum\limits_{j = 1}^N u_j^*(\boldsymbol{z})  p_i(\boldsymbol{x}_j) = \sum\limits_{i = 1}^q  g_i(\boldsymbol{z} ) .
	$$ 
	
	Then $f^*(\boldsymbol{z})$ can be rewritten as 
	\begin{equation}
		f^*(\boldsymbol{z}) = \sum\limits_{i = 1}^q  g_i(\boldsymbol{z}) f(\xi_i) +\left\langle f, \sum\limits_{j = 1}^N u_j^*(\boldsymbol{z})\kappa(\cdot,\boldsymbol{x}_j)- \sum\limits_{i= 1}^q p_i(\boldsymbol{z})\kappa(\cdot,\xi_i)\right\rangle_\mathcal{H}. 
	\end{equation}
	
	In the other hand, the functions $g_1(x), g_2(x),\dots,g_q(x)$ are also linearly independent, then 
	\begin{equation}
		f(\boldsymbol{z}) = \sum\limits_{i = 1}^q f(\xi_i) g_i(\boldsymbol{z}) +\left\langle f,\kappa(\cdot,\boldsymbol{z})-\sum\limits_{i= 1}^q g_i(\boldsymbol{z})\kappa(\cdot,\xi_i) \right\rangle_\mathcal{H}.
	\end{equation}
	
	Thus 
	$$
	\begin{aligned}
		&\vert f(\boldsymbol{z}) - f^*(\boldsymbol{z})\vert^2 =\left\vert \left \langle f,\kappa(\cdot,\boldsymbol{z}) - \sum\limits_{j = 1}^N u_j^*(\boldsymbol{z})\kappa(\cdot,\boldsymbol{x}_j) \right \rangle_\mathcal{H} \right\vert^2 \\
		&\le \Vert f\Vert_\mathcal{H}^2 \cdot \left\vert \left \langle \kappa(\cdot,\boldsymbol{z}) - \sum\limits_{j = 1}^N u_j^*(\boldsymbol{z})\kappa(\cdot,\boldsymbol{x}_j), \kappa(\cdot,\boldsymbol{z}) - \sum\limits_{j = 1}^N u_j^*(\boldsymbol{z})\kappa(\cdot,\boldsymbol{x}_j) \right \rangle_\mathcal{H} \right\vert\\
		&= \Vert f\Vert_\mathcal{H}^2 \cdot \left(\kappa(\boldsymbol{z},\boldsymbol{z})- K_X(\boldsymbol{z}) K_{XX}^{-1} K_X(\boldsymbol{z})\right) = C\cdot Var(f^*(\boldsymbol{z})), 
	\end{aligned}
	$$
	where $C =\Vert f\Vert_\mathcal{H}^2$. 
\end{proof}

\end{document}
